Table~\ref{tab:main} gives the main results at default weights and Fig.~\ref{fig:effort} the control effort.

\begin{table*}[t]\centering\caption{Main results at default weights, mean $\pm$ sample std over 5 seeds (20 random starts each). Bold: best, underline: second best. Tied values are marked by the table generator in row order, so ties (e.g.\ cart-pole success \rhval{main/ilqr-energy-ours/cartpole/success_rate/mean:2}) carry no ranking information.}\label{tab:main}
\begin{tabular}{llccccc}
\toprule
Method & Task & success\_rate $\uparrow$ & success\_rate\_far $\uparrow$ & control\_effort $\downarrow$ & time\_to\_upright $\downarrow$ & $n$ \\
\midrule
EnergyShaping-LQR & pendulum & \textbf{0.98 $\pm$ 0.04} & \textbf{0.96 $\pm$ 0.10} & \textbf{10.78 $\pm$ 1.66} & \textbf{2.25 $\pm$ 0.27} & 5 \\
iLQR-Quad & pendulum & 0.80 $\pm$ 0.08 & 0.52 $\pm$ 0.16 & 17.24 $\pm$ 3.68 & 2.75 $\pm$ 0.30 & 5 \\
iLQR-Quad+Energy & pendulum & 0.95 $\pm$ 0.04 & 0.88 $\pm$ 0.08 & \underline{12.23 $\pm$ 2.20} & 2.33 $\pm$ 0.20 & 5 \\
\textbf{iLQR-Energy (ours)} & pendulum & \underline{0.95 $\pm$ 0.04} & \underline{0.88 $\pm$ 0.08} & 12.24 $\pm$ 2.22 & \underline{2.33 $\pm$ 0.20} & 5 \\
\midrule
EnergyShaping-LQR & cartpole & 0.97 $\pm$ 0.04 & 0.93 $\pm$ 0.10 & \underline{54.91 $\pm$ 3.66} & 2.70 $\pm$ 0.12 & 5 \\
iLQR-Quad & cartpole & \underline{1.00 $\pm$ 0.00} & \underline{1.00 $\pm$ 0.00} & 74.77 $\pm$ 3.48 & \underline{1.71 $\pm$ 0.07} & 5 \\
iLQR-Quad+Energy & cartpole & 1.00 $\pm$ 0.00 & 1.00 $\pm$ 0.00 & 69.52 $\pm$ 2.44 & \textbf{1.66 $\pm$ 0.07} & 5 \\
\textbf{iLQR-Energy (ours)} & cartpole & \textbf{1.00 $\pm$ 0.00} & \textbf{1.00 $\pm$ 0.00} & \textbf{42.03 $\pm$ 6.03} & 1.79 $\pm$ 0.06 & 5 \\
\bottomrule
\end{tabular}
\end{table*}

\begin{figure}[t]\centering\includegraphics[width=\columnwidth]{main_effort.pdf}
\caption{Control effort per system and task at default weights (lower is better). Bars: mean over 5 seeds; error bars: $\pm$ one sample std over seeds. Quad+E: iLQR-Quad+Energy; E-shaping: EnergyShaping-LQR.}\label{fig:effort}\end{figure}

\textbf{H1 (success over quadratic): not supported.} On the pendulum iLQR-Energy succeeds more often than iLQR-Quad (\rhval{main/ilqr-energy-ours/pendulum/success_rate/mean:2} vs \rhval{main/ilqr-quad/pendulum/success_rate/mean:2}; paired $p=\rhval{cmp/main/ilqr-quad/pendulum/success_rate/paired_p:3}$, $n=5$ seeds), and more so on far starts (\rhval{main/ilqr-energy-ours/pendulum/success_rate_far/mean:2} vs \rhval{main/ilqr-quad/pendulum/success_rate_far/mean:2}). On the cart-pole both reach \rhval{main/ilqr-energy-ours/cartpole/success_rate/mean:2}, a ceiling, so there is no difference to detect. By our pre-stated criterion (higher on \emph{both} tasks) H1 is not supported. The pendulum gap also holds at default weights only (Sec.~\ref{sec:ablations}).

\textbf{H2 (lower effort): supported at default weights, smaller after re-weighting.} By the registered test (main group, control effort) iLQR-Energy has lower effort than iLQR-Quad on both tasks (pendulum \rhval{main/ilqr-energy-ours/pendulum/control_effort/mean:2} vs \rhval{main/ilqr-quad/pendulum/control_effort/mean:2}, paired $p=\rhval{cmp/main/ilqr-quad/pendulum/control_effort/paired_p:3}$; cart-pole \rhval{main/ilqr-energy-ours/cartpole/control_effort/mean:2} vs \rhval{main/ilqr-quad/cartpole/control_effort/mean:2}, $p=\rhval{cmp/main/ilqr-quad/cartpole/control_effort/paired_p:4}$). The size of the cart-pole gap is mostly a product of the default $Q$. In the sweep (Table~\ref{tab:sweep}, seeds 0--2) iLQR-Quad at $q=0.1$ keeps success \rhval{sweep_qscale/ilqr-quad@qscale=0.1/cartpole/success_rate/mean:2} and a time to upright of \rhval{sweep_qscale/ilqr-quad@qscale=0.1/cartpole/time_to_upright/mean:2}\,s with effort \rhval{sweep_qscale/ilqr-quad@qscale=0.1/cartpole/control_effort/mean:1}, against \rhval{sweep_we/ilqr-energy-ours@we=1/cartpole/control_effort/mean:1} for iLQR-Energy at $w_E=1$. The gap therefore shrinks from \rhval{cmp/pair_q1/ilqr-energy-ours/cartpole/control_effort/delta:1} to \rhval{cmp/pair_q0.1/ilqr-energy-ours/cartpole/control_effort/delta:1} on those seeds, where it is not clearly distinguishable (Table~\ref{tab:paired}). With 5 seeds the tuned comparison gives \rhval{tuned_all/ilqr-energy-ours/cartpole/control_effort/mean:1} vs \rhval{tuned_all/ilqr-quad/cartpole/control_effort/mean:1} (Table~\ref{tab:tuned}). On the pendulum the direction persists across the whole $q$ range (\rhval{sweep_we/ilqr-energy-ours@we=1/pendulum/control_effort/mean:1} vs \rhval{sweep_qscale/ilqr-quad@qscale=100/pendulum/control_effort/mean:1}--\rhval{sweep_qscale/ilqr-quad@qscale=1/pendulum/control_effort/mean:1}). Effort is measured over episodes that include failures, and iLQR-Quad reaches the goal nominally earlier on the cart-pole at default weights (\rhval{main/ilqr-quad/cartpole/time_to_upright/mean:2} vs \rhval{main/ilqr-energy-ours/cartpole/time_to_upright/mean:2}\,s; not distinguishable at $n=5$); a speed/effort trade-off is a possible explanation, not a demonstrated one.

\textbf{H4 (versus energy shaping): mixed.} On the pendulum EnergyShaping-LQR has the highest success (\rhval{main/energyshaping-lqr/pendulum/success_rate/mean:2}) and lowest effort. Its difference from iLQR-Energy is not distinguishable with 5 seeds for success (\rhval{main/energyshaping-lqr/pendulum/success_rate/mean:2} vs \rhval{main/ilqr-energy-ours/pendulum/success_rate/mean:2}) or time to upright (\rhval{main/energyshaping-lqr/pendulum/time_to_upright/mean:2} vs \rhval{main/ilqr-energy-ours/pendulum/time_to_upright/mean:2}\,s), but it uses less effort (\rhval{main/energyshaping-lqr/pendulum/control_effort/mean:2} vs \rhval{main/ilqr-energy-ours/pendulum/control_effort/mean:2}, paired $p=\rhval{cmp/main/energyshaping-lqr/pendulum/control_effort/paired_p:3}$, lower in all 5 seeds). On the cart-pole iLQR-Energy is better on success (\rhval{main/ilqr-energy-ours/cartpole/success_rate/mean:2} vs \rhval{main/energyshaping-lqr/cartpole/success_rate/mean:2}, not distinguishable), effort (\rhval{main/ilqr-energy-ours/cartpole/control_effort/mean:2} vs \rhval{main/energyshaping-lqr/cartpole/control_effort/mean:2}) and time to upright (\rhval{main/ilqr-energy-ours/cartpole/time_to_upright/mean:2} vs \rhval{main/energyshaping-lqr/cartpole/time_to_upright/mean:2}\,s). By the registered test (success and time to upright) H4 is supported on the cart-pole and not refuted on the pendulum, where energy shaping is nominally better on both and better on effort. The energy-shaping gains were not tuned.

\textbf{Energy alone vs.\ combined.} On the pendulum iLQR-Energy and iLQR-Quad+Energy are equal in success (\rhval{main/ilqr-energy-ours/pendulum/success_rate/mean:2}) and within \rhval{cmp/main/ilqr-quad+energy/pendulum/control_effort/delta:2} in effort (\rhval{main/ilqr-energy-ours/pendulum/control_effort/mean:2} vs \rhval{main/ilqr-quad+energy/pendulum/control_effort/mean:2}). On the cart-pole they differ: the combined cost has a faster upright time (\rhval{main/ilqr-quad+energy/cartpole/time_to_upright/mean:2} vs \rhval{main/ilqr-energy-ours/cartpole/time_to_upright/mean:2}\,s, paired $p=\rhval{cmp/main/ilqr-quad+energy/cartpole/time_to_upright/paired_p:3}$) but higher effort (\rhval{main/ilqr-quad+energy/cartpole/control_effort/mean:2} vs \rhval{main/ilqr-energy-ours/cartpole/control_effort/mean:2}, paired $p=\rhval{cmp/main/ilqr-quad+energy/cartpole/control_effort/paired_p:4}$), closer to iLQR-Quad (\rhval{main/ilqr-quad/cartpole/control_effort/mean:2}). To see which term carries the combined cost we logged, along its closed-loop trajectories, the mean of the two state terms $\tfrac12(E-E^{*})^2$ and $\tfrac12 e^\top Q e$ (Table~\ref{tab:diag}). The energy term is \rhval{diag_terms/ilqr-quad+energy/pendulum/term_ratio/mean:2} times the quadratic term on the pendulum and \rhval{diag_terms/ilqr-quad+energy/cartpole/term_ratio/mean:2} times on the cart-pole (mean over seeds of the per-run ratio), where only the light pole's energy is counted. This is consistent with the combined cost behaving like the energy cost on the pendulum and more like the quadratic cost on the cart-pole. It is a measurement of magnitudes, not of gradients, so the terms' sizes may explain the pattern but do not prove it.

\begin{table}[t]\centering\caption{Closed-loop mean of the two state-cost terms under iLQR-Quad+Energy at default weights (seeds 0--2, all time steps); ratio is energy over quadratic, computed per run and then averaged over the seeds.}\label{tab:diag}
\begin{tabular}{lrrr}\toprule
Task & energy term & quad.\ term & ratio \\\midrule
pendulum & \rhval{diag_terms/ilqr-quad+energy/pendulum/energy_term/mean:2} & \rhval{diag_terms/ilqr-quad+energy/pendulum/quad_term/mean:2} & \rhval{diag_terms/ilqr-quad+energy/pendulum/term_ratio/mean:2} \\
cart-pole & \rhval{diag_terms/ilqr-quad+energy/cartpole/energy_term/mean:2} & \rhval{diag_terms/ilqr-quad+energy/cartpole/quad_term/mean:2} & \rhval{diag_terms/ilqr-quad+energy/cartpole/term_ratio/mean:2} \\
\bottomrule
\end{tabular}
\end{table}

\textbf{Where the method fails.} The pendulum failures of iLQR-Energy are more frequent among far starts (\rhval{main/ilqr-energy-ours/pendulum/success_rate_far/mean:2} far vs \rhval{main/ilqr-energy-ours/pendulum/success_rate/mean:2} overall), the regime the cost was meant to help; energy shaping does better there (\rhval{main/energyshaping-lqr/pendulum/success_rate_far/mean:2}). With a small weight the energy cost loses its advantage on the pendulum: at $w_E=0.01$ success is \rhval{sweep_we/ilqr-energy-ours@we=0.01/pendulum/success_rate/mean:2} and effort \rhval{sweep_we/ilqr-energy-ours@we=0.01/pendulum/control_effort/mean:1}, the level of iLQR-Quad (Table~\ref{tab:sweep}). On the cart-pole a large weight costs effort (\rhval{sweep_we/ilqr-energy-ours@we=100/cartpole/control_effort/mean:1} at $w_E=100$).

\textbf{What this does not show.} It does not show that the energy cost is more robust under model error or noise (none was tested), nor that it generalises beyond two low-dimensional plants.
